\chapter{Regressões Aparentemente Não-Relacionadas}
\label{cap:sur}
% ── índice ──────────────────────────────────────────────────────
\index{SUR|textbf}
\index{seemingly unrelated regressions|see{SUR}}
\index{Zellner, estimador de|textbf}
\index{sur@\texttt{sur()}|textbf}
\index{eficiência, ganho de}
\index{correlação contemporânea}

% ─────────────────────────────────────────────────────────────────────────────
\section{O estimador de Zellner}

Quando se estima um sistema de $M$ equações com os mesmos (ou diferentes)
regressores e os erros são correlacionados entre equações, o OLS equação
a equação é ineficiente. O estimador \textbf{SUR} (\emph{Seemingly Unrelated
Regressions}) de Zellner (1962) explora essa correlação via GLS conjunto.

Sistema de $M$ equações:
\[
  y_m = X_m\,\beta_m + \varepsilon_m, \quad m = 1,\ldots,M
\]

Empilhando: $Y = X\beta + \varepsilon$, com $E[\varepsilon\varepsilon'] = \Sigma \otimes I_n$,
onde $\Sigma$ é a matriz de covariância entre-equações ($M \times M$).

O estimador GLS conjunto:
\[
  \hat\beta_{\mathrm{SUR}} = (X'(\Sigma^{-1}\otimes I)X)^{-1}\,X'(\Sigma^{-1}\otimes I)\,Y
\]

\textbf{Ganho de eficiência}: SUR $\geq$ OLS quando $\Sigma$ tem elementos
fora da diagonal não nulos — i.e., quando os erros das equações são
correlacionados. Se $\Sigma$ for diagonal ou se os regressores forem idênticos,
SUR $=$ OLS.

% ─────────────────────────────────────────────────────────────────────────────
\section{Verificação por DGP}

DGP com dois erros fortemente correlacionados ($\rho = 0{,}7$):

\[
  y_1 = 1 + 2\,x_1 + \varepsilon_1 \qquad y_2 = 0{,}5 + 1{,}5\,x_2 + \varepsilon_2
\]
\[
  \mathrm{Cov}(\varepsilon_1, \varepsilon_2) = 0{,}7\,\sigma_1\,\sigma_2
\]

\begin{lstlisting}[caption={DGP SUR e comparação com OLS equação a equação}]
seed(42)
generate df x1 = rnormal()
generate df x2 = rnormal()
generate df e1 = (rnormal() - 0.5) * 1.7321
generate df e2 = 0.7 * e1 + 0.7141 * (rnormal() - 0.5) * 1.7321
generate df y1 = 1.0 + 2.0 * x1 + e1
generate df y2 = 0.5 + 1.5 * x2 + e2

let m_ols1 = ols(y1 ~ x1, df)
let m_ols2 = ols(y2 ~ x2, df)
esttab(m_ols1, m_ols2)

let m_sur = sur(df, y1 ~ x1, y2 ~ x2)
display m_sur
\end{lstlisting}

\subsection*{OLS equação a equação}

\begin{lstlisting}[language={}, basicstyle=\ttfamily\small,
  frame=none, numbers=none, backgroundcolor=\color{codbg},
  xleftmargin=0pt, xrightmargin=0pt, breaklines=false]
                            (1)              (2)
                            OLS              OLS
────────────────────────────────────────────────
x1                    2.1956***
                       (0.0788)
x2                                     1.5097***
                                        (0.0779)
Constant              0.9050***        0.4879***
────────────────────────────────────────────────
N                           500              500
R²                       0.6091           0.4300
\end{lstlisting}

\subsection*{SUR}

\begin{lstlisting}[language={}, basicstyle=\ttfamily\small,
  frame=none, numbers=none, backgroundcolor=\color{codbg},
  xleftmargin=0pt, xrightmargin=0pt, breaklines=false]
══════════════════════════════════════════════════════════════════════════════
               Seemingly Unrelated Regressions (SUR)
         Cross-Equation Error Correlation (Σ):
  [  0.2463   0.1689 ]
  [  0.1689   0.2419 ]
────────────────────────────────────────────────
 Equation: y1  —  const=0.9513  x1=2.1003***   R²=0.6080
 Equation: y2  —  const=0.4591  x2=1.5653***   R²=0.4294
══════════════════════════════════════════════════════════════════════════════
\end{lstlisting}

A correlação estimada entre erros é $\hat\Sigma_{12}/\sqrt{\hat\Sigma_{11}\hat\Sigma_{22}} = 0{,}169/\sqrt{0{,}246 \times 0{,}242} \approx 0{,}69$, próximo do verdadeiro $0{,}7$. O SUR ganha eficiência via corte nos SEs: $\hat\sigma_{\mathrm{SUR}}(x_1) < \hat\sigma_{\mathrm{OLS}}(x_1)$.

% ─────────────────────────────────────────────────────────────────────────────
\section{Sintaxe}

\begin{lstlisting}
sur(df, y1 ~ x1 + x2, y2 ~ x3 + x4)
sur(df, y1 ~ x1, y2 ~ x2, y3 ~ x3)        // 3 equações
\end{lstlisting}

O primeiro argumento é o DataFrame; os demais são fórmulas (uma por equação).

\begin{nota}
  Se todas as equações tiverem os mesmos regressores, SUR = OLS e não há
  ganho de eficiência. O ganho é máximo quando os regressores são ortogonais
  entre equações e a correlação entre erros é alta.
\end{nota}

% ─────────────────────────────────────────────────────────────────────────────
\section{Validação empírica}

O dataset de investimento \texttt{grunfeld} do Wooldridge é usado no caso de
validação \texttt{sur_grunfeld} em \texttt{validation/cases/}. O caso estima um
sistema SUR de duas equações no \hayashi{} e compara os coeficientes com
implementações de referência em R e Python. Execute \lstinline{hay validate} para
ver o estado atual.
